\documentclass[11pt, a4paper]{article}
\usepackage[UTF8]{ctex}

\usepackage[margin=2.5cm]{geometry}
\usepackage{setspace}
\setstretch{1.3}

\usepackage{amsmath,amssymb,amsfonts}
\usepackage{mathtools}
\usepackage{booktabs}
\usepackage{tabularx}
\usepackage{array}
\usepackage{enumitem}
\usepackage{graphicx}
\usepackage{xcolor}
\usepackage[most]{tcolorbox}
\usepackage{hyperref}
\usepackage{microtype}
\usepackage{fancyhdr}
\usepackage{titlesec}
\usepackage{abstract}

% ---- 颜色定义 ----
\definecolor{myblue}{RGB}{25, 70, 150}
\definecolor{myred}{RGB}{160, 30, 30}
\definecolor{mygray}{RGB}{248, 248, 250}
\definecolor{mybluegray}{RGB}{230, 238, 250}

% ---- 页眉页脚 ----
\pagestyle{fancy}
\fancyhf{}
\fancyhead[L]{\small\textit{Powell (2022): RL and Stochastic Optimization}}
\fancyhead[R]{\small\textit{贾云凯}}
\fancyfoot[C]{\small\thepage}
\renewcommand{\headrulewidth}{0.4pt}

% ---- 标题格式 ----
\titleformat{\section}{\large\bfseries\color{myblue}}{{\color{myblue}\thesection.}}{0.5em}{}[\vspace{-0.3em}\rule{\textwidth}{0.8pt}]
\titleformat{\subsection}{\normalsize\bfseries\color{myblue!80!black}}{{\thesubsection}}{0.5em}{}
\titleformat{\paragraph}[runin]{\normalsize\bfseries}{}{0em}{}[.\quad]

% ---- tcolorbox 样式 ----
\tcbset{
  keybox/.style={
    colback=mybluegray,
    colframe=myblue,
    coltitle=white,
    fonttitle=\bfseries\small,
    title=#1,
    sharp corners,
    boxrule=1pt,
    top=5pt, bottom=5pt, left=8pt, right=8pt,
  },
  formulabox/.style={
    colback=blue!4,
    colframe=myblue,
    sharp corners,
    boxrule=0.7pt,
    top=5pt, bottom=5pt, left=8pt, right=8pt,
  },
  scenariobox/.style={
    colback=orange!5,
    colframe=orange!70!black,
    fonttitle=\bfseries\small,
    coltitle=white,
    sharp corners,
    boxrule=0.7pt,
    top=5pt, bottom=5pt, left=8pt, right=8pt,
  },
  warnbox/.style={
    colback=red!4,
    colframe=myred,
    fonttitle=\bfseries\small,
    coltitle=white,
    title=#1,
    sharp corners,
    boxrule=0.7pt,
    top=5pt, bottom=5pt, left=8pt, right=8pt,
  }
}

\hypersetup{colorlinks=true, linkcolor=myblue, citecolor=myblue, urlcolor=myblue}

\begin{document}

% ========== 封面 ==========
\begin{titlepage}
\centering
\vspace*{1.5cm}
{\color{myblue}\rule{\textwidth}{2pt}}
\vspace{0.8cm}

{\LARGE\bfseries\color{myblue}
\textit{Reinforcement Learning and Stochastic Optimization}\\[0.4em]
A Unified Framework for Sequential Decisions}

\vspace{0.4cm}
{\large Warren B. Powell \quad Princeton University \quad Wiley, 2022}

\vspace{0.6cm}
{\color{myblue}\rule{\textwidth}{2pt}}

\vspace{1.2cm}
{\huge\bfseries 书籍导读与第一章精读}

\vspace{0.5cm}
{\large 贾云凯}

\vspace{2cm}
\begin{tcolorbox}[formulabox, width=0.85\textwidth]
\centering
\textbf{全书核心目标函数：}
\[
  \max_{\pi = (f \in \mathcal{F},\; \theta \in \Theta_f)}\ 
  \mathbb{E}\!\left\{\sum_{t=0}^{T} C(S_t,\ X^\pi(S_t)) \,\middle|\, S_0\right\}
\]
\textit{寻找最优策略 $\pi$，使期望累积回报最大——这是贯穿全书的统一框架。}
\end{tcolorbox}

\vfill
{\normalsize 整理日期：\today}
\end{titlepage}

% ========== 摘要 ==========
\begin{abstract}
本文对 Powell（2022）进行系统性介绍。全书核心论断是：强化学习、随机规划、马尔可夫决策过程、最优控制、多臂赌博机等15个研究社区，本质上研究的是同一类问题——\textbf{顺序决策问题}。作者提出"先建模，后求解（Model first, then solve）"的方法论，以\textbf{五元组通用建模框架}统一刻画所有问题，并将所有策略方法归纳为\textbf{四大类型}：策略函数近似（PFA）、代价函数近似（CFA）、价值函数近似（VFA）与直接前瞻近似（DLA）。本文展开第一章核心内容，辅以具体公式和应用场景，供读者快速把握全书方法论框架。
\end{abstract}

\tableofcontents
\newpage

%% ===================================================================
\section{为什么需要"统一框架"？}
%% ===================================================================

\subsection{学术背景与核心痛点}

顺序决策问题（Sequential Decision Problems）广泛存在于工程、经济、医学、物流与能源等领域。然而，长期以来，各研究社区在\textbf{符号体系、建模范式与算法偏好}上高度割裂：

\begin{itemize}[leftmargin=2em, itemsep=2pt]
  \item \textbf{计算机科学}（强化学习、MDP）：偏爱离散动作，倾向于价值函数近似与Q-learning；
  \item \textbf{控制工程}（最优控制、MPC）：偏爱连续状态-动作，倾向于前瞻优化与Hamilton-Jacobi方程；
  \item \textbf{运筹学}（随机规划、仿真优化）：偏爱向量决策与情景树（scenario tree）；
  \item \textbf{统计与机器学习}（多臂赌博机、主动学习）：偏爱纯学习问题与Bayesian方法。
\end{itemize}

同一问题在不同社区有不同名字、不同符号、不同算法，导致\textbf{知识交叉与方法迁移极为困难}。本书的贡献，正是在15个社区之上建立统一的语言、建模框架与策略分类体系。

\begin{tcolorbox}[keybox={作者的核心声明（可证伪的通用性命题）}]
任何顺序决策问题的\textbf{求解策略}必然属于如下四类之一（或其混合）：
\[
  \pi \in \{\underbrace{\text{PFA}}_{\text{策略函数近似}},\quad 
            \underbrace{\text{CFA}}_{\text{代价函数近似}},\quad 
            \underbrace{\text{VFA}}_{\text{价值函数近似}},\quad 
            \underbrace{\text{DLA}}_{\text{直接前瞻近似}}\}
\]
这一分类\textbf{覆盖学术文献与工程实践中所有已知方法}。
\end{tcolorbox}

\subsection{全书章节结构}

\begin{center}
\begin{tabular}{clp{9.5cm}}
\toprule
\textbf{部分} & \textbf{章节} & \textbf{内容要点} \\
\midrule
Part I & 第1--4章 & 导论、通用建模框架、在线学习、随机搜索入门 \\
Part II & 第5--7章 & 状态无关问题：梯度型/无梯度随机搜索、步长策略 \\
Part III & 第8--11章 & 状态依赖问题：建模方法、不确定性建模、策略设计 \\
Part IV & 第12--13章 & 策略搜索类：PFA（策略函数近似）、CFA（代价函数近似） \\
Part V & 第14--19章 & 前瞻类策略：精确DP、近似ADP/RL、直接前瞻（MCTS等） \\
Part VI & 第20章 & 多智能体系统与学习 \\
\bottomrule
\end{tabular}
\end{center}

%% ===================================================================
\section{第一章精读：顺序决策问题的本质}
%% ===================================================================

\subsection{问题的最简定义}

\paragraph{顺序决策问题（Sequential Decision Problem）} 其本质为以下序列：
\[
  \underbrace{x_0}_{\text{决策}},\; 
  \underbrace{W_1}_{\text{信息}},\; 
  \underbrace{x_1}_{\text{决策}},\; 
  \underbrace{W_2}_{\text{信息}},\; \ldots
\]
形式化地，完整样本路径写成：
\begin{equation}
  (S_0,\; x_0,\; W_1,\; S_1,\; x_1,\; W_2,\; \ldots,\; S_t,\; x_t,\; W_{t+1},\; \ldots,\; S_T)
  \label{eq:sample_path}
\end{equation}
其中 $S_t$ 为\textbf{状态变量}，$x_t$ 为\textbf{决策变量}，$W_{t+1}$ 为决策后才能观察到的\textbf{外生随机信息}。

\begin{tcolorbox}[scenariobox, title={场景一：库存管理（最经典的顺序决策问题）}]
某仓库每期需决定订货量 $x_t \geq 0$，随机需求 $\hat{D}_{t+1}$ 在下期才能观察到。库存演化方程为：
\begin{equation}
  R_{t+1} = \max\{0,\; R_t + x_t - \hat{D}_{t+1}\}
  \label{eq:inventory}
\end{equation}
这正是序列 \eqref{eq:sample_path} 的具体实例：每期决策订货量 $\to$ 随机需求到达 $\to$ 库存更新，循环往复。其中不确定性来源于需求 $\hat{D}_{t+1}$，而 $x_t$ 的选择直接影响下一期状态 $R_{t+1}$。
\end{tcolorbox}

\subsection{五元组通用建模框架}

Powell 将任何顺序决策问题统一为如下\textbf{五元组建模框架}：

\begin{tcolorbox}[keybox={通用建模框架（Universal Modeling Framework）}]
\begin{enumerate}[leftmargin=1.5em, label=\textbf{(\arabic*)}, itemsep=3pt]
  \item \textbf{状态变量 $S_t$：} 包含做决策和建模所需的全部信息（且仅含必要信息）；
  \item \textbf{决策变量 $x_t$：} 可为二元、离散、连续标量或向量；决策规则记为 $x_t = X^\pi(S_t)$；
  \item \textbf{外生信息 $W_{t+1}$：} 做出决策 $x_t$ 后才到来的外部随机信息；
  \item \textbf{转移函数 $S^M(S_t, x_t, W_{t+1})$：} 描述状态如何从 $t$ 演化到 $t+1$；
  \item \textbf{目标函数：} 最大化期望累积回报（见式 \eqref{eq:objective}）。
\end{enumerate}
\end{tcolorbox}

\noindent\textbf{标准目标函数：}
\begin{equation}
  \max_{\pi}\; \mathbb{E}\!\left\{
    \sum_{t=0}^{T} C(S_t,\; X^\pi(S_t)) \;\middle|\; S_0
  \right\}
  \label{eq:objective}
\end{equation}
其中 $C(S_t, x_t)$ 为单期贡献（收益/回报/负代价），$\pi$ 为策略（从状态到决策的映射函数），期望 $\mathbb{E}$ 对 $W_1, \ldots, W_T$ 的所有随机性取均值，$S_0$ 为初始状态（可含参数不确定性分布）。

当策略 $\pi$ 含可调参数 $\theta \in \Theta_f$ 与函数类 $f \in \mathcal{F}$ 时，目标函数推广为：
\begin{equation}
  \max_{\pi = (f \in \mathcal{F},\; \theta \in \Theta_f)}\; \mathbb{E}\!\left\{
    \sum_{t=0}^{T} C(S_t,\; X^\pi(S_t)) \;\middle|\; S_0
  \right\}
  \label{eq:objective_full}
\end{equation}

\begin{tcolorbox}[scenariobox, title={场景一（续）——五元组建模库存问题}]
\begin{enumerate}[leftmargin=1.5em, label=\textbullet, itemsep=1pt]
  \item \textbf{状态：} $S_t = R_t$（当前库存水平）；
  \item \textbf{决策：} $x_t \geq 0$（订货量）；策略 $X^\pi(S_t)$ 待设计；
  \item \textbf{外生信息：} $W_{t+1} = \hat{D}_{t+1}$（随机需求）；
  \item \textbf{转移：} $R_{t+1} = \max\{0,\; R_t + x_t - \hat{D}_{t+1}\}$；
  \item \textbf{目标：} $\max_\pi \mathbb{E}\!\left\{\sum_{t=0}^{T} [p\min\{R_t+x_t,\hat{D}_{t+1}\} - cx_t] \mid S_0\right\}$，\\
        其中 $p$ 为销售单价，$c$ 为订货单价。
\end{enumerate}
\end{tcolorbox}

\subsection{状态变量的三层结构}

状态变量 $S_t$ 由三类信息构成：$S_t = (R_t,\; I_t,\; B_t)$

\begin{center}
\begin{tabular}{p{2.5cm}p{2cm}p{9cm}}
\toprule
\textbf{类型} & \textbf{符号} & \textbf{含义与示例} \\
\midrule
物理状态 & $R_t$ & 库存水平、无人机位置、资产持仓量、能量存储水平 \\
\addlinespace
已知信息 & $I_t$ & 当前已知价格、确定性合同参数、日历信息 \\
\addlinespace
信念/置信状态 & $B_t$ & 对未知参数的概率分布估计（Bayesian 后验） \\
\bottomrule
\end{tabular}
\end{center}

\noindent\textbf{注：} $B_t$ 是区别于经典MDP的重要元素。信念状态使模型能自然处理\textbf{主动学习（Active Learning）}——即通过选择信息收集策略来更新对未知参数的认知。

\subsection{方法论核心："先建模，后求解"}

\begin{tcolorbox}[keybox={核心方法论：先建模后求解}]
\textbf{第一步（建模）：} 完整写出五元组：$S_t$、$x_t \in \mathcal{X}_t$、$W_{t+1}$、$S^M(\cdot)$、目标函数 \eqref{eq:objective}。\\[4pt]
\textbf{第二步（求解）：} 完成建模后，从四类策略中选择适合的方法（或其混合）。
\end{tcolorbox}

\noindent 这与确定性优化完全平行：先写出线性规划/整数规划模型，再选择求解器。顺序决策中，"求解器"变成了"策略设计"。其根本区别在于：
\begin{itemize}[leftmargin=2em, itemsep=1pt]
  \item 确定性优化寻找\textbf{实数向量} $(x_0, x_1, \ldots, x_T)$；
  \item 顺序决策寻找\textbf{函数（策略）} $\pi$：$S_t \mapsto x_t$。
\end{itemize}

%% ===================================================================
\section{四类策略：通用分类体系}
%% ===================================================================

所有策略分为两大方向，每个方向两类，共四类：

\begin{equation}
  \underbrace{
    \overbrace{\text{PFA}}^{\text{策略函数近似}},\quad
    \overbrace{\text{CFA}}^{\text{代价函数近似}}
  }_{\text{策略搜索（Policy Search）}}\qquad\qquad
  \underbrace{
    \overbrace{\text{VFA}}^{\text{价值函数近似}},\quad
    \overbrace{\text{DLA}}^{\text{直接前瞻近似}}
  }_{\text{前瞻近似（Lookahead）}}
\end{equation}

\subsection{策略函数近似（PFA）}

分析函数直接将状态映射到决策，\textbf{无需嵌套优化}：
\begin{equation}
  X^{\text{PFA}}(S_t \,|\, \theta) = \theta_0 + \sum_{f=1}^{F} \theta_f \phi_f(S_t)
  \label{eq:pfa_linear}
\end{equation}
或非线性形式（如神经网络、查找表、分段线性函数）。参数 $\theta$ 通过仿真或在线调参确定。

\begin{tcolorbox}[scenariobox, title={PFA 应用——库存问题的订货上限策略}]
\begin{equation}
  X^{\text{Inv}}(S_t \,|\, \theta) = \begin{cases}
    \theta^{\max} - R_t, & \text{若 } R_t < \theta^{\min}, \\
    0, & \text{否则}.
  \end{cases}
  \label{eq:inventory_pfa}
\end{equation}
参数 $\theta = (\theta^{\min}, \theta^{\max})$ 通过求解下式确定：
\[
  \max_{\theta}\; \mathbb{E}\!\left\{\sum_{t=0}^{T} C(S_t,\; X^{\text{Inv}}(S_t|\theta),\; W_{t+1})\right\}
\]
这是实践中\textbf{最广泛使用}的策略类型，简洁但调参本身具有挑战性。
\end{tcolorbox}

\subsection{代价函数近似（CFA）}

参数化（嵌套）优化模型，内部含 $\arg\max$ 子问题：
\begin{equation}
  X^{\text{CFA}}(S_n \,|\, \theta) = \arg\max_{x \in \mathcal{X}}\!\left(\bar{\mu}_x^n + \theta\, \bar{\sigma}_x^n\right)
  \label{eq:cfa_ucb}
\end{equation}
式 \eqref{eq:cfa_ucb} 称为\textbf{区间估计（Interval Estimation / UCB）}，其中 $\bar{\mu}_x^n$ 为均值估计，$\bar{\sigma}_x^n$ 为标准差，$\theta \geq 0$ 调节探索-利用权衡强度。

\begin{tcolorbox}[scenariobox, title={CFA 应用——在线广告选择（多臂赌博机）}]
广告集合 $\mathcal{X} = \{x_1, \ldots, x_M\}$，每次展示后观察到点击（$W_{n+1}=1$）或不点击（$W_{n+1}=0$）。

UCB 策略以 \eqref{eq:cfa_ucb} 选择下一个展示的广告，其五元组建模为：
\begin{itemize}[leftmargin=1.5em, itemsep=1pt]
  \item $S_n = \{(\bar{\mu}_x^n, \bar{\sigma}_x^n)\}_{x \in \mathcal{X}}$（信念状态 $B_n$）；
  \item $x_n \in \mathcal{X}$：展示的广告；
  \item $W_{n+1} \in \{0,1\}$：点击观测；
  \item 转移函数：Bayesian 更新 $(\bar{\mu}_x^n, \bar{\sigma}_x^n)$；
  \item 目标：$\max_\theta \mathbb{E}\!\left\{\sum_{n=0}^{N} W_{n+1}(x_n)\right\}$（最大化累积点击量）。
\end{itemize}
\end{tcolorbox}

CFA 的核心特征是\textbf{内嵌优化问题}：$x$ 可以是简单的离散选择，也可以是大规模整数规划（如航班排班、能源调度）。

\subsection{价值函数近似（VFA）}

基于 Bellman 方程的前瞻策略：
\begin{equation}
  V_t(S_t) = \max_{x_t \in \mathcal{X}_t}\!\left(C(S_t, x_t) + \mathbb{E}_{W_{t+1}}\!\{V_{t+1}(S_{t+1}) \mid S_t\}\right)
  \label{eq:bellman}
\end{equation}
最优策略为：
\begin{equation}
  X_t^{\text{VFA}}(S_t) = \arg\max_{x_t \in \mathcal{X}_t}\!\left(C(S_t, x_t) + \mathbb{E}_{W_{t+1}}\!\{\bar{V}_{t+1}(S_{t+1}) \mid S_t\}\right)
  \label{eq:vfa_policy}
\end{equation}
其中 $\bar{V}_{t+1}$ 为近似价值函数（由机器学习方法，如神经网络、线性模型、查找表，从仿真数据中估计）。强化学习中的 Q-learning、TD($\lambda$)、Actor-Critic 均属此类。

\begin{tcolorbox}[warnbox={重要警示（对 VFA/RL 研究者的反思）}]
\begin{enumerate}[leftmargin=1.5em, label=\textbullet, itemsep=1pt]
  \item Bellman 方程 \eqref{eq:bellman} 在绝大多数现实问题中\textbf{不可直接精确求解}（维度诅咒）；
  \item 近似价值函数 $\bar{V}$ 的拟合往往\textbf{不稳定}，尤其在高维连续状态空间；
  \item 作者明确指出：VFA 是``a powerful methodology for a \emph{surprisingly narrow} set of decision problems''；
  \item \textbf{Bellman 方程并非最优策略的唯一路径}——PFA 和 DLA 同样可以包含最优策略。
\end{enumerate}
\end{tcolorbox}

\subsection{直接前瞻近似（DLA）}

在未来 $H$ 步内进行（近似）规划，执行当前步最优决策：
\begin{equation}
  X_t^{\text{DLA}}(S_t) = \arg\max_{x_t}\!\left[C(S_t, x_t) + \max_{\pi'} \mathbb{E}\!\left\{\sum_{t'=t+1}^{t+H} C(S_{t'}, X^{\pi'}(S_{t'})) \;\middle|\; S_t, x_t\right\}\right]
  \label{eq:dla}
\end{equation}

\begin{center}
\begin{tabular}{lp{11cm}}
\toprule
\textbf{DLA 变体} & \textbf{说明} \\
\midrule
确定性前瞻 & 用点预测 $\hat{W}_{t+1}, \ldots, \hat{W}_{t+H}$ 求解未来 $H$ 期确定性优化（如导航系统最短路径、模型预测控制MPC） \\
\addlinespace
随机前瞻 & 对未来随机性显式建模，求解嵌套随机优化（两阶段随机规划等） \\
\addlinespace
MCTS & 蒙特卡洛树搜索，通过随机仿真估计树节点的期望价值 \\
\bottomrule
\end{tabular}
\end{center}

\subsection{四类策略对比总结}

\begin{center}
\begin{tabular}{p{1.6cm}p{3.2cm}p{3.5cm}p{4cm}}
\toprule
\textbf{类型} & \textbf{核心机制} & \textbf{代表方法} & \textbf{典型适用场景} \\
\midrule
PFA & 解析函数映射 $S_t \to x_t$ & 神经网络策略、线性决策律 & 低维连续动作、工程落地最广 \\
\addlinespace
CFA & 参数化嵌套优化（给已有优化模型加可学习参数） & UCB/置信区间加权、参数化目标优化、鲁棒修正项 & 多臂赌博机、高维资源分配 \\
\addlinespace
VFA & Bellman 方程近似 & Q-learning、TD($\lambda$)、ADP、Actor-Critic & 库存优化、储能调度、时序决策、游戏AI \\
\addlinespace
DLA & 有限步前瞻优化（滚动时域 / 树搜索） & MCTS、MPC、多阶段随机优化（有限滚动时域） & 需精确建模未来影响的决策 \\
\bottomrule
\end{tabular}
\end{center}

\begin{tcolorbox}[scenariobox, title={四类策略的通俗理解}]
\textbf{PFA（策略函数近似）—— “端到端直接映射，实时无优化”}  
最简单、工程落地最广的范式：给定当前状态 $S_t$，通过固定或可学习参数的显式函数直接输出决策 $x_t$，决策时刻不再求解任何优化问题。  
\textbf{例子}：库存管理中经典的“(s, S)阈值策略”——看到库存低于 $\theta^{\min}$ 就订货补到 $\theta^{\max}$。公式定义为：
\[
X^{\text{PFA}}(S_t|\theta) = 
\begin{cases}
\theta^{\max} - R_t, & R_t < \theta^{\min} \\
0, & \text{否则}
\end{cases}
\]
参数 $\theta = (\theta^{\min},\theta^{\max})$ 可通过仿真、强化学习或梯度方法离线标定，是工业时序调度中最常见的落地形式。

\textbf{CFA（代价函数近似）—— “原生优化 + 可学习微调旋钮”}  
依托成熟的优化模型（规划、调度、资源分配），**不改变模型结构**，仅引入可学习参数 $\theta$ 来修正目标函数、约束惩罚或置信权重，实现探索-利用平衡或鲁棒性自适应。  
\textbf{例子}：在线广告选择（多臂赌博机）中的UCB策略：
\[
X^{\text{CFA}}(S_n|\theta) = \arg\max_x \left( \bar{r}_x + \theta \cdot \text{CI}(x) \right)
\]
其中 $\theta$ 越大越倾向探索新广告。本质是把已有优化框架保留，只对“代价”进行参数化微调。

\textbf{VFA（价值函数近似）—— “先拟合未来收益，再用Bellman判优”}  
不直接给出决策，而是先训练一个价值函数 $\bar{V}_t(S_t)$（或Q函数），近似估计“如果以后一直最优，未来还能赚多少钱”；再通过Bellman最优性准则选择当前动作。  
\textbf{例子}：储能调度中，当前电价低时要不要充电？VFA会评估“充电后未来价值增加 $\bar{V}_{t+1}(E_{t+1})$”，从而决定是否值得。这是ADP和强化学习在能源、库存等工业时序决策中的标准用法。Q-learning、TD($\lambda$)、Actor-Critic均属此类。

\textbf{DLA（直接前瞻近似）—— “显式滚动仿真，算完H步只走第一步”}  
不拟合价值函数，而是直接基于系统模型或随机采样，求解未来有限时域 $H$ 内的全路径优化；只执行第一步决策，下一期重新预测、重新规划。  
\textbf{例子}：模型预测控制（MPC）——预测未来24小时电价，解一个24期优化问题，只执行第1期的动作；明天再滚动重算。MCTS（蒙特卡洛树搜索）和有限滚动时域随机规划也属于此类。
\end{tcolorbox}

\noindent\textbf{小结}  
PFA = 显式决策映射（快、易落地）；  
CFA = 优化模型参数微调（适配资源分配 / 赌博机）；  
VFA = Bellman价值拟合（长期收益最优）；  
DLA = 有限时域前瞻仿真（高精度强建模）。  

工程前沿几乎都采用**混合架构**：如MPC（DLA）内嵌VFA修正终端价值、资源调度（CFA）用PFA输出快速基线解，从而大幅兼顾精度与计算效率。

\begin{tcolorbox}[formulabox]
\textbf{核心权衡：简洁性 vs. 计算代价}
\[
  \underbrace{\text{PFA, CFA}}_{\text{简洁，实践广泛}} \implies \underbrace{\text{含可调参数 } \theta}_{\text{调参本身是难题！}}
\qquad\qquad
  \underbrace{\text{VFA, DLA}}_{\text{学术热点}} \implies \underbrace{\text{计算代价高}}_{\text{可行性约束强}}
\]
\end{tcolorbox}

%% ===================================================================
\section{从确定性优化到随机优化：迁移路径}
\label{sec:det2sto}
%% ===================================================================

考虑如下确定性多期规划问题：
\begin{align}
  \max_{x_0, \ldots, x_T}\quad & \sum_{t=0}^{T} p_t x_t \label{eq:det_obj} \\
  \text{s.t.}\quad & A_t x_t = R_t,\quad x_t \geq 0 \\
  & R_{t+1} = B_t x_t + \hat{R}_{t+1}
\end{align}
当价格 $p_t$ 随机变化、供给 $\hat{R}_{t+1}$ 为随机变量时，通过以下\textbf{四步迁移}转化为随机顺序决策：

\begin{center}
\begin{tabular}{lll}
\toprule
\textbf{修改项} & \textbf{确定性} & \textbf{随机化后} \\
\midrule
决策对象 & 实数向量 $x_t$ & 函数（策略） $X^\pi(S_t)$ \\
贡献函数 & $C_t(x_t)$ & $C(S_t, x_t)$（依赖状态，含随机价格 $p_t$） \\
优化方式 & $\max_{x_0,\ldots,x_T}$ & $\max_\pi$（搜索策略空间） \\
目标值 & $\sum_t p_t x_t$ & $\mathbb{E}\!\left\{\sum_t C(S_t, X^\pi(S_t))\right\}$ \\
\bottomrule
\end{tabular}
\end{center}

随机化后的目标函数：
\begin{equation}
  \max_{\pi}\; \mathbb{E}\!\left\{
    \sum_{t=0}^{T} C(S_t,\; X^\pi(S_t)) \;\middle|\; S_0
  \right\}
  \label{eq:sto_final}
\end{equation}

\noindent\textbf{注意——约束的随机化不是简单替换：} 原全局预算约束 $\sum_{t=0}^T x_t \leq B$ 需改写为\textbf{递归约束}：
\begin{align}
  x_t &\leq B - R_t^{\text{budget}}, \label{eq:rec_constr} \\
  R_{t+1}^{\text{budget}} &= R_t^{\text{budget}} + x_t. \label{eq:rec_update}
\end{align}
$R_t^{\text{budget}}$ 纳入状态变量，策略 $X^\pi(S_t)$ 在每期必须满足约束 \eqref{eq:rec_constr}。

%% ===================================================================
\section{应用场景展开}
%% ===================================================================

\subsection{场景二：资产出售——最优停止（PFA + VFA 对比）}

令 $p_t$ 为资产价格，$R_t \in \{0,1\}$ 表示是否持有资产，决策 $x_t=1$ 代表出售。

\noindent\textbf{价格演化模型（均值回归）：}
\begin{equation}
  p_{t+1} = p_t + \theta(p_t - \bar{p}) + \varepsilon_{t+1},\quad \varepsilon_{t+1} \sim \mathcal{N}(0, \sigma_\varepsilon^2)
\end{equation}

\noindent\textbf{PFA 策略（阈值型）：}
\begin{equation}
  X^{\text{PFA}}(S_t \,|\, \rho) = \mathbf{1}\!\left[p_t \geq \bar{p}_t + \rho,\; R_t = 1\right],
  \quad \bar{p}_t = 0.9\bar{p}_{t-1} + 0.1 p_t
\end{equation}

\noindent\textbf{VFA 策略（Bellman 递推）：}
\begin{equation}
  V_t(p_t, R_t) = \max_{x_t \in \{0,1\}}\!\left[p_t x_t R_t + (1-x_t)\mathbb{E}_{p_{t+1}}\!\{V_{t+1}(p_{t+1}, R_t)\}\right]
\end{equation}
因状态空间连续，精确求解不可行，需参数化近似 $\bar{V}_{t+1}$（如线性模型或高斯过程）。

\subsection{场景三：能量存储调度——四类策略对比实验}

设储能系统当前电量 $E_t$，可充电 $u_t^+ \geq 0$ 或放电 $u_t^- \geq 0$，实时电价 $p_t$ 随机波动。状态 $S_t = (E_t, p_t)$，转移方程：
\[
  E_{t+1} = E_t + \eta u_t^+ - u_t^-/\eta,\quad 0 \leq E_{t+1} \leq E^{\max}
\]
其中 $\eta \in (0,1]$ 为充放电效率。四类策略的实例化：

\begin{center}
\begin{tabular}{lp{11cm}}
\toprule
\textbf{策略类} & \textbf{具体形式} \\
\midrule
PFA & 当 $p_t < \theta_1$（低价）时充电，当 $p_t > \theta_2$（高价）时放电；参数 $\theta = (\theta_1, \theta_2)$ 待调 \\
\addlinespace
CFA & 修改目标函数为 $\max [p_t(u_t^- - u_t^+) + \theta \cdot \bar{V}_{t+1}(E_{t+1})]$，通过参数 $\theta$ 调节对未来价值的估计权重 \\
\addlinespace
VFA & $\arg\max_{u_t}\!\left[C(S_t, u_t) + \bar{V}_{t+1}(E_{t+1})\right]$，用分段线性凸函数近似 $\bar{V}$（保证凸性） \\
\addlinespace
DLA & 以未来 $H$ 期预测价格 $\hat{p}_{t+1}, \ldots, \hat{p}_{t+H}$ 求解确定性前瞻优化（即模型预测控制MPC） \\
\bottomrule
\end{tabular}
\end{center}

\noindent\textbf{关键实验结论（Powell 书中第11章）：} 通过调整价格波动率、储能容量、预测精度等参数，可以使四类策略\textbf{各自表现最佳}。这验证了"先建模，后从四类策略中选择"的必要性——\textbf{没有任何一类策略对所有问题普遍最优}。

%% ===================================================================
\section{机器学习与顺序决策的形式联系}
%% ===================================================================

\subsection{形式对比}

\begin{center}
\begin{tabular}{ll}
\toprule
\textbf{统计/机器学习} & \textbf{顺序决策（随机优化）} \\
\midrule
\textbf{批量学习：} $\min_\theta \frac{1}{N}\sum_{n=1}^N (y^n - f(x^n|\theta))^2$ & \textbf{样本均值近似（SAA）：} $\min_{x \in \mathcal{X}} \frac{1}{N}\sum_{n=1}^N F(x, W(\omega^n))$ \\
\addlinespace
\textbf{在线学习：} $\min_\theta \mathbb{E}(Y - f(X|\theta))^2$ & \textbf{随机搜索：} $\min_\theta \mathbb{E}\,F(X, W)$ \\
\addlinespace
\textbf{函数类搜索：} $\min_{f \in \mathcal{F}, \theta \in \Theta_f} \mathbb{E}(Y - f(X|\theta))^2$ & \textbf{策略搜索：} $\min_\pi \mathbb{E}\!\sum_{t=0}^T C(S_t, X^\pi(S_t))$ \\
\bottomrule
\end{tabular}
\end{center}

\noindent\textbf{核心区别：} 机器学习有给定的训练集 $(x^n, y^n)$；顺序决策以转移函数 $S^M$ 和外生信息过程 $W_{t+1}$ 替代——通过与环境交互\textbf{生成}数据，而非使用预先给定的数据集。

\subsection{在线学习在四类策略中的角色}

在线学习贯穿全书，具体体现在：
\begin{enumerate}[leftmargin=2em, label=\textbf{(\arabic*)}, itemsep=2pt]
  \item \textbf{估计期望：} 用随机梯度/蒙特卡洛近似 $\mathbb{E}_W F(x, W)$（Part II 核心）；
  \item \textbf{调整策略参数 $\theta$：} PFA/CFA 的在线调参；
  \item \textbf{近似价值函数 $\bar{V}_t(S_t)$：} VFA/ADP/RL 的核心学习任务；
  \item \textbf{学习转移模型：} 当 $S^M$ 未知时的 model-free 方法（第9章）。
\end{enumerate}

%% ===================================================================
\section{阅读建议}
%% ===================================================================

\subsection{章节优先级}

\begin{center}
\begin{tabular}{p{2cm}p{3.5cm}p{8cm}}
\toprule
\textbf{优先级} & \textbf{章节} & \textbf{推荐理由} \\
\midrule
\textbf{必读} & 第1、9、11章 & 通用框架建立；完整建模方法；策略设计导论（含实验对比） \\
\addlinespace
\textbf{强烈推荐} & 第7章 & 无梯度随机搜索：四类策略在纯学习问题中的第一次完整展示 \\
\addlinespace
\textbf{按需阅读} & 第12--13章 & PFA/CFA 详述：实践最常用，学术关注度低，值得深入 \\
\addlinespace
\textbf{深度学习} & 第14--17章 & 精确DP、近似ADP、Q-learning、策略梯度；RL 从业者必读 \\
\addlinespace
\textbf{进阶} & 第19章 & 直接前瞻（MCTS、两阶段随机规划）；需较强数学基础 \\
\bottomrule
\end{tabular}
\end{center}

\subsection{核心批判性思考问题}

\begin{enumerate}[leftmargin=2em, label=\textbf{Q\arabic*.}, itemsep=4pt]
  \item 在你当前的研究问题中，状态变量 $S_t = (R_t, I_t, B_t)$ 三层结构各包含哪些要素？信念状态 $B_t$ 如何通过转移函数更新？

  \item 你所在研究领域的主流方法属于四类策略中的哪一类？是否存在其他类别被长期忽视？

  \item 如何将你熟悉的一个确定性优化模型，按第 \ref{sec:det2sto} 节的四步迁移路径，转化为式 \eqref{eq:sto_final} 的随机顺序决策形式？

  \item 你研究问题中的 VFA 类方法（强化学习）是否真的优于 PFA/CFA？能否设计受控实验加以验证？

  \item ``调参的困难''（"The price of simplicity is tunable parameters, and tuning is hard!"）在你的问题中如何体现？是否可以量化调参代价？
\end{enumerate}

\subsection{符号体系速查表}

\begin{center}
\begin{tabular}{ll|ll}
\toprule
\textbf{符号} & \textbf{含义} & \textbf{符号} & \textbf{含义} \\
\midrule
$S_t$ & 状态变量（含 $R_t, I_t, B_t$） & $X^\pi(S_t)$ & 策略/决策函数 \\
$x_t$ & 决策变量 & $\pi$ & 策略（函数族） \\
$W_{t+1}$ & 外生随机信息 & $\theta$ & 策略参数 \\
$S^M(\cdot)$ & 状态转移函数 & $C(S_t, x_t)$ & 单期贡献函数 \\
$V_t(S_t)$ & 价值函数 & $\mathbb{E}\{\cdot\}$ & 期望算子 \\
$\bar{V}_t(S_t)$ & 近似价值函数 & $T$ & 决策时域长度 \\
$\mathcal{X}_t$ & 决策可行集 & $\mathcal{F}$ & 策略函数类 \\
\bottomrule
\end{tabular}
\end{center}

%% ===================================================================
\section{总结}
%% ===================================================================

本书最重要的三点贡献：

\begin{enumerate}[leftmargin=2em, label=\textbf{\arabic*.}, itemsep=4pt]
  \item \textbf{统一建模语言：} 五元组框架 $(S_t, x_t, W_{t+1}, S^M, \text{目标函数})$ 可描述任意顺序决策问题，打破15个研究社区的符号壁垒，使跨社区方法比较成为可能。

  \item \textbf{通用策略分类：} PFA / CFA / VFA / DLA 四类策略覆盖所有已知决策方法。这是一个可被挑战的\textbf{通用性声明}，为不同社区提供了共同比较语言。

  \item \textbf{方法论转变：} 从``为问题找算法''到``先建模，后从工具箱中选策略''——更接近工程实践，更有利于教学与研究的系统化。
\end{enumerate}

\noindent 对博士研究生而言，本书最大价值或许不在某一具体算法，而在于提供了一个\textbf{元框架（meta-framework）}：面对新的决策问题时，能够系统性地追问："我的状态变量是什么？外生信息的分布模型是什么？四类策略中，哪一类最适合这个问题的结构？"

\begin{tcolorbox}[formulabox]
\begin{center}
\textbf{全书统一目标函数——值得反复思考的核心：}
\end{center}
\begin{equation}
  \boxed{
  \max_{\pi = (f \in \mathcal{F},\; \theta \in \Theta_f)}\; \mathbb{E}\!\left\{
    \sum_{t=0}^{T} C(S_t,\; X^\pi(S_t)) \;\middle|\; S_0
  \right\}
  }
  \label{eq:final_unified}
\end{equation}
\begin{center}
策略 $\pi$ 涵盖 PFA、CFA、VFA、DLA 及其混合形式。\\
这是 Powell（2022）对顺序决策领域40年研究的核心凝练。
\end{center}
\end{tcolorbox}

%% 参考文献
\begin{thebibliography}{9}
\bibitem{powell2022}
Powell, W.~B. (2022).
\textit{Reinforcement Learning and Stochastic Optimization: A Unified Framework for Sequential Decisions}.
John Wiley \& Sons.

\bibitem{powell2019}
Powell, W.~B. (2019).
A unified framework for stochastic optimization.
\textit{European Journal of Operational Research}, 275(3), 795--821.

\bibitem{powell2014}
Powell, W.~B. (2014).
Clearing the jungle of stochastic optimization.
\textit{INFORMS Tutorials in Operations Research: Bridging Data and Decisions}, 109--137.

\bibitem{sutton2018}
Sutton, R.~S. \& Barto, A.~G. (2018).
\textit{Reinforcement Learning: An Introduction}, 2nd ed. MIT Press.

\bibitem{powell2011}
Powell, W.~B. (2011).
\textit{Approximate Dynamic Programming: Solving the Curses of Dimensionality}, 2nd ed. John Wiley \& Sons.
\end{thebibliography}

\end{document}